\documentclass[11pt,letterpaper]{article}
\usepackage[margin=1in]{geometry}
\usepackage{setspace}
\usepackage{enumitem}
\usepackage{array}
\usepackage{booktabs}
\usepackage{tabularx}
\usepackage{graphicx}
\usepackage{titlesec}
\usepackage{fancyhdr}
\usepackage[T1]{fontenc}
\usepackage[utf8]{inputenc}
\usepackage{lmodern}

\setstretch{1.03}
\pagestyle{empty}
\setlength{\parindent}{15pt}
\setlength{\parskip}{4pt}
\setlist{nosep,leftmargin=0.24in}
\titlespacing*{\section}{0pt}{0.35em}{0.1em}
\titlespacing*{\subsection}{0pt}{0.25em}{0.08em}
\titleformat{\section}{\normalfont\bfseries\large}{}{0pt}{}
\titleformat{\subsection}{\normalfont\bfseries}{}{0pt}{}

\begin{document}
\begin{center}
{\Large\bfseries Facilities, Equipment and Other Resources}\\[0.05in]
{\normalsize Mostafa Rezaali}\\
{\normalsize NSF AGS-PRF Application}
\end{center}
\vspace{-0.08in}
\section*{Host Institution}
The host institution is University of Florida, Department of Geography. The Department of Geography provides an appropriate environment for a fellowship centered on climate extremes, hydroclimatology, geospatial analysis, environmental prediction, and AI-enabled climate-risk modeling. The department provides access to scholarly mentoring, seminars, research collaboration, and professional-development opportunities relevant to the proposed work.

\section*{Sponsoring Scientist and Intellectual Environment}
The sponsoring scientist is David Keellings, University of Florida. The sponsor's expertise in heat waves, climate extremes, and applied climatology directly supports the proposed research. The intellectual environment will help ensure that model development is evaluated against physical understanding of atmospheric variability, land-atmosphere coupling, and regional climate risk.

\section*{Computational Resources}
The project will use Python/PyTorch software environments, geospatial and climate-data processing tools, and GPU-enabled computing resources available through the host research environment and institutional computing pathways. The existing codebase includes scripts for model training, cross-validation, hindcast export, bootstrap statistics, baseline comparison, temporal anomaly correlation maps, sample maps, and probabilistic verification. These resources are adequate for training deterministic and probabilistic graph models, storing derived hindcast products, and generating reproducible diagnostics.

\section*{Data Resources}
The project will use publicly available and/or institutionally accessible climate datasets, including PRISM daily maximum temperature fields, ERA5 atmospheric variables, teleconnection indices, topography, land masks, seasonal radiation, and land-surface variables. Raw third-party datasets will be obtained through their official providers. Derived products needed to reproduce figures and tables will be documented and preserved according to the Data Management and Sharing Plan.

\section*{Software and Analysis Environment}
The analysis environment includes Python, scientific computing libraries, geospatial processing tools, NetCDF/GRIB workflows, visualization packages, and version-control practices. The project will emphasize reproducible workflow design, clear metadata, and portable scripts for model training and verification. No specialized laboratory equipment is required beyond computing, storage, and standard research software.

\section*{Professional Development and Dissemination Resources}
The host environment will support manuscript preparation, conference presentations, student mentoring, seminars, and development of teaching material related to trustworthy AI for climate extremes. These resources will help the Fellow build an independent research identity in probabilistic subseasonal prediction and climate-risk analytics.
\end{document}
